\chapter{刚体运动学}

\section{刚体上各点的速度和加速度}

\begin{problem}{5.1.1}
揉茶机的揉桶有三个曲柄支承，曲柄的支座为 $A, B, C$ ，支轴为 $a, b, c$ ，且 $A, B, C$ 与 $a, b, c$ 为两等边三角形，各曲柄长 $r$ ，均以不变的角速度 $\omega$ 绕其支座转动，如图5.1所示，求揉桶中心 $O$ 点的速率和加速度的大小。
\begin{center}
\begin{tikzpicture}[>=Stealth, line width=0.8pt]
  \def\R{1.7}
  \def\Rear{1.95}
  \def\rcrank{0.72}
  \def\crankang{45}

  \coordinate (O) at (0,0);
  \coordinate (c) at (90:\Rear);
  \coordinate (a) at (210:\Rear);
  \coordinate (b) at (330:\Rear);

  \coordinate (C) at ($(c)-(\crankang:\rcrank)$);
  \coordinate (A) at ($(a)-(\crankang:\rcrank)$);
  \coordinate (B) at ($(b)-(\crankang:\rcrank)$);

  % Drum outline with three smooth rounded lobes at 90, 210, 330 deg
  \draw
    (102:\R) arc (102:198:\R)
    .. controls (202:1.9) and (204:2.12) .. (210:2.12)
    .. controls (216:2.12) and (218:1.9) .. (222:\R)
    arc (222:318:\R)
    .. controls (322:1.9) and (324:2.12) .. (330:2.12)
    .. controls (336:2.12) and (338:1.9) .. (342:\R)
    arc (342:438:\R)
    .. controls (82:1.9) and (84:2.12) .. (90:2.12)
    .. controls (96:2.12) and (98:1.9) .. (102:\R);

  % Drum center O
  \fill (O) circle (1.4pt);
  \node[right, inner sep=3pt] at (O) {$O$};

  % Helper-like repetition for the three cranks and fixed supports
  \foreach \P/\p/\labP/\labp/\posP/\posp in {
    A/a/A/a/left/above right,
    B/b//b/left/above,
    C/c/C/c/left/right%
  } {
    % Fixed support triangle and hatched ground
    \draw ($(\P)+(-0.18,-0.28)$) -- (\P) -- ($(\P)+(0.18,-0.28)$);
    \draw ($(\P)+(-0.3,-0.28)$) -- ($(\P)+(0.3,-0.28)$);
    \foreach \x in {-0.22,-0.11,0,0.11,0.22} {
      \draw ($(\P)+(\x,-0.28)$) -- ($(\P)+(\x-0.08,-0.38)$);
    }

    % Crank bar (two parallel edges + pivot circles)
    \draw ($(\P)+(\crankang+90:0.06)$) -- ($(\p)+(\crankang+90:0.06)$);
    \draw ($(\P)+(\crankang-90:0.06)$) -- ($(\p)+(\crankang-90:0.06)$);
    \draw[fill=white] (\P) circle (2.4pt);
    \draw[fill=white] (\p) circle (2.4pt);

    % Crank length label r and angular velocity arrow omega
    \node[above left, inner sep=1pt] at ($(\P)!0.65!(\p)+(\crankang+90:0.06)$) {$r$};
    \draw[->] ($(\P)+(105:0.30)$) arc (105:-15:0.30) node[right, inner sep=1pt] {$\omega$};

    % Pivot labels
    \node[\posP, inner sep=4pt] at (\P) {$\labP$};
    \node[\posp, inner sep=3.5pt] at (\p) {$\labp$};
  }
  \node[below right, inner sep=1pt] at ($(B)+(0.24, -0.26)$) {$B$};
\end{tikzpicture}

\end{center}
\end{problem}
\clearpage


\begin{problem}{5.1.2}
半径为 r 的半圆形凸轮以匀速 u 沿水平线向右运动, 带动直杆 AB 绕 A 点转动, 当 $\theta = 30^{\circ}$ 时, A 点与凸轮的中心 O 点正好在同一铅直线上(见图 5.2). 求此时 AB 杆与凸轮相接触的一点 C 的速度与加速度的大小.
\begin{center}
\begin{tikzpicture}[>=Stealth, line width=0.8pt]
  \def\r{1.75}
  \def\hw{0.07}
  \pgfmathsetmacro{\H}{2*(\r+\hw)} % A=(0,H) so left edge of rod is tangent at theta=30 deg

  \coordinate (O) at (0,0);
  \coordinate (A) at (0,\H);
  \coordinate (C) at (30:\r);

  % Ground line and hatches
  \draw (-2.3, 0) -- (2.3, 0);
  \foreach \x in {-2.15,-1.95,...,2.15} {
    \draw (\x, 0) -- (\x-0.12, -0.16);
  }
  \node[below, inner sep=3.5pt] at (0, -0.14) {$O$};

  % Semicircle cam
  \draw (\r, 0) arc (0:180:\r);

  % Radius r arrow
  \draw[->] (O) -- (140:\r) node[midway, below left, inner sep=1pt] {$r$};

  % Vertical dashed line OA
  \draw[dashed] (O) -- (A);

  % Rod AB (centerline starts at A, angle -60 deg, half-width hw)
  \coordinate (Bend) at ($(A)+(-60:3.72)$);
  \draw[fill=white]
    ($(A)+(-150:\hw)$) -- ($(Bend)+(-150:\hw)$) -- ($(Bend)+(30:\hw)$) -- ($(A)+(30:\hw)$) -- cycle;

  % Ceiling support at A
  \draw ($(A)+(-0.26,0.42)$) -- (A) -- ($(A)+(0.26,0.42)$);
  \draw (-0.42, \H+0.42) -- (0.42, \H+0.42);
  \foreach \x in {-0.35,-0.20,-0.05,0.10,0.25,0.40} {
    \draw (\x, \H+0.42) -- (\x+0.12, \H+0.55);
  }
  \draw[fill=white] (A) circle (2.6pt);

  % Angle theta
  \draw ($(A)+(-90:0.75)$) arc (-90:-60:0.75);
  \node at ($(A)+(-75:1.05)$) {$\theta$};

  % Labels A, C, B
  \node[right, inner sep=4pt] at (A) {$A$};
  \node[below left, inner sep=2pt] at (C) {$C$};
  \node[right, inner sep=2pt] at (Bend) {$B$};
\end{tikzpicture}

\end{center}
\end{problem}
\clearpage


\begin{problem}{5.1.3}
图 5.4 中画的是蒸汽机或内燃机中用的曲柄、滑块连杆机构。曲柄 OA 以等角速度 $\omega$ 绕 O 点转动，通过连杆 AB 带动活塞 B 在滑槽中左右运动。若 OA = R, AB = l, t = 0 时 $\varphi = 0$ 。求 t 时活塞 B 的速度。
\begin{center}
\begin{tikzpicture}[>=Stealth, line width=0.8pt]
  \def\R{1.25}
  \def\L{3.1}
  \def\phiang{48}
  \pgfmathsetmacro{\Ax}{\R*cos(\phiang)}
  \pgfmathsetmacro{\Ay}{\R*sin(\phiang)}
  \pgfmathsetmacro{\Bx}{\Ax + sqrt(\L*\L - \Ay*\Ay)}
  \pgfmathsetmacro{\thetaang}{asin(\Ay/\L)}

  \coordinate (O) at (0,0);
  \coordinate (A) at (\Ax, \Ay);
  \coordinate (B) at (\Bx, 0);

  % Crank circle (dashed)
  \draw[dashed] (O) circle (\R);

  % Horizontal centerline
  \draw[dashed] (O) -- (\Bx-0.35, 0);

  % Fixed support at O
  \draw (-0.18,-0.28) -- (O) -- (0.18,-0.28);
  \draw (-0.32,-0.28) -- (0.32,-0.28);
  \foreach \x in {-0.24,-0.12,0,0.12,0.24} {
    \draw (\x,-0.28) -- (\x-0.08,-0.38);
  }

  % Guide channel walls for piston B
  \draw (\Bx-0.45, 0.22) -- (\Bx+0.75, 0.22);
  \draw (\Bx-0.45, -0.22) -- (\Bx+0.75, -0.22);
  \foreach \x in {-0.4,-0.22,-0.04,0.14,0.32,0.50,0.68} {
    \draw (\Bx+\x, 0.22) -- (\Bx+\x+0.10, 0.34);
    \draw (\Bx+\x, -0.22) -- (\Bx+\x-0.10, -0.34);
  }

  % Piston slider block at B
  \draw[fill=white] (\Bx-0.36, 0.16) -- (\Bx+0.22, 0.16) -- (\Bx+0.22, -0.16) -- (\Bx-0.36, -0.16);
  \draw[dashed] (\Bx-0.36, 0.16) -- (\Bx-0.36, -0.16);

  % Connecting rod AB (rounded tip into piston)
  \pgfmathsetmacro{\rodang}{-\thetaang}
  \draw[fill=white]
    ($(A)+(\rodang+90:0.06)$) -- ($(B)+(\rodang+90:0.06)$)
    arc (\rodang+90:\rodang-90:0.06) -- ($(A)+(\rodang-90:0.06)$) -- cycle;

  % Crank OA
  \draw[fill=white]
    ($(O)+(\phiang+90:0.07)$) -- ($(A)+(\phiang+90:0.07)$) --
    ($(A)+(\phiang-90:0.07)$) -- ($(O)+(\phiang-90:0.07)$) -- cycle;
  \draw[fill=white] (O) circle (2.5pt);
  \draw[fill=white] (A) circle (2.5pt);

  % Angles phi and theta
  \draw (0.45, 0) arc (0:\phiang:0.45);
  \node at ({0.5*\phiang}:0.66) {$\varphi$};
  \draw ($(B)+(180:0.75)$) arc (180:{180-\thetaang}:0.75);
  \node at ($(B)+({180-0.5*\thetaang}:1.05)$) {$\theta$};

  % Angular velocity arrow omega
  \draw[->] (-25:0.35) arc (-25:120:0.35) node[above left, inner sep=1pt] {$\omega$};

  % Labels
  \node[left, inner sep=3.5pt] at (O) {$O$};
  \node[above left, inner sep=1pt] at ($(O)!0.55!(A)+(\phiang+90:0.08)$) {$R$};
  \node[above, inner sep=3.5pt] at (A) {$A$};
  \node[above right, inner sep=1pt] at ($(A)!0.48!(B)$) {$l$};
  \node[above, inner sep=2pt] at (\Bx+0.15, 0.36) {$B$};
\end{tikzpicture}

\end{center}
\end{problem}
\clearpage


\begin{problem}{5.1.4}
曲柄 $OA$ 以匀角速 $\omega_0 = 10\mathrm{s}^{-1}$ 顺时针方向转动(见图5.5). 求当 $\theta = 30^{\circ}$ 时 $BC$ 杆的角速度 $\omega$ 、角加速度 $\beta$ 及销钉 $A$ 在 $BC$ 杆导槽中的相对加速度 $a_r$ .
\begin{center}
\begin{tikzpicture}[>=Stealth, line width=0.8pt]
  \def\R{1.6}
  \coordinate (O) at (0,0);
  \coordinate (C) at (\R, 0);
  \coordinate (A) at (120:\R);

  % Dashed circle centered at O
  \draw[dashed] (O) circle (\R);

  % Horizontal dashed baseline through O and C
  \draw[dashed] (-1.2, 0) -- (C);

  % Radius label 20cm
  \draw (O) -- (205:\R) node[midway, below, inner sep=4pt] {$20\,\mathrm{cm}$};

  % Fixed supports at O and C
  \foreach \P in {O, C} {
    \draw ($(\P)+(-0.18,-0.28)$) -- (\P) -- ($(\P)+(0.18,-0.28)$);
    \draw ($(\P)+(-0.32,-0.28)$) -- ($(\P)+(0.32,-0.28)$);
    \foreach \x in {-0.24,-0.12,0,0.12,0.24} {
      \draw ($(\P)+(\x,-0.28)$) -- ($(\P)+(\x-0.08,-0.38)$);
    }
  }

  % Crank OA (angle 120 deg, extends slightly past A)
  \coordinate (OAtop) at (120:{\R+0.18});
  \draw[fill=white]
    ($(O)+(210:0.12)$) arc (210:390:0.12) --
    ($(OAtop)+(30:0.12)$) arc (30:210:0.12) -- cycle;

  % Slotted link BC along direction 150 deg from C
  \pgfmathsetmacro{\LenAC}{sqrt(3)*\R}
  \coordinate (B) at ($(C)+(150:{\LenAC+0.9})$);
  \draw[fill=white]
    ($(C)+(60:0.16)$) -- ($(B)+(60:0.16)$) arc (60:240:0.16) --
    ($(C)+(240:0.16)$) arc (240:420:0.16) -- cycle;

  % Inner slot of BC
  \coordinate (S1) at ($(C)+(150:0.35)$);
  \coordinate (S2) at ($(C)+(150:{\LenAC+0.55})$);
  \draw
    ($(S1)+(60:0.055)$) -- ($(S2)+(60:0.055)$) arc (60:240:0.055) --
    ($(S1)+(240:0.055)$) arc (240:420:0.055) -- cycle;

  % Hidden part of OA behind BC (subtle dashed edge)
  \draw[dashed] ($(A)+(30:0.12)$) -- ($(O)!0.75!(A)+(30:0.12)$);

  % Pivots and pin A
  \draw[fill=white] (O) circle (2.5pt);
  \draw[fill=white] (C) circle (2.5pt);
  \fill (A) circle (2.0pt);

  % Angle 2theta at O and theta at C
  \draw (-0.52, 0) arc (180:120:0.52);
  \node at (150:0.78) {$2\theta$};
  \draw ($(C)+(180:0.65)$) arc (180:150:0.65);
  \node at ($(C)+(165:0.95)$) {$\theta$};

  % Labels
  \node[above right, inner sep=3.5pt] at ($(O)+(0.03, 0.02)$) {$O$};
  \node[above right, inner sep=3.5pt] at ($(C)+(0.03, 0.02)$) {$C$};
  \node[above right, inner sep=2pt] at (120:{\R+0.26}) {$A$};
  \node[above left, inner sep=2pt] at (B) {$B$};
\end{tikzpicture}

\end{center}
\end{problem}
\clearpage


\begin{problem}{5.1.5}
摇床急用机构的曲柄以不变转速 $n = 90\mathrm{r / min}$ 绕定轴 $O$ 转动，通过滑块 $A$ 带动扇形齿轮绕 $O_{1}$ 轴转动，从而带动齿条 $B$ 做铅垂方向振动。设 $OO_{1}$ 在同一水平线上，曲柄 $OA = 76\mathrm{mm}$ ，其他尺寸如图5.6所示。

求 $\theta = 30^{\circ}$ 时齿条 $B$ 的速度和加速度.
\begin{center}
\begin{tikzpicture}[>=Stealth, line width=0.8pt]
  \def\Rgear{1.52}
  \def\Dist{1.83}
  \def\rcrank{0.76}
  \def\cranktheta{50}

  \coordinate (O1) at (0,0);
  \coordinate (O) at (\Dist, 0);
  \coordinate (A) at ({\Dist - \rcrank*cos(\cranktheta)}, {-\rcrank*sin(\cranktheta)});
  \pgfmathsetmacro{\armang}{atan2(-\rcrank*sin(\cranktheta), \Dist - \rcrank*cos(\cranktheta))}

  % Vertical rack B on the left (right face at x = -Rgear)
  \def\xrack{-1.52}
  \draw (\xrack-0.42, -1.65) -- (\xrack-0.42, 1.55);
  \draw[dashed] (\xrack-0.12, -1.65) -- (\xrack-0.12, 1.55);
  \draw (\xrack, -1.65) -- (\xrack, 1.55);
  \node[left, inner sep=3pt] at (\xrack-0.42, 0.1) {$B$};

  % Rack guide blocks (top and bottom, left and right)
  \foreach \y in {-1.35, 1.25} {
    % Left block
    \draw (\xrack-0.66, \y-0.16) rectangle (\xrack-0.52, \y+0.16);
    \foreach \dy in {-0.12, 0, 0.12} {
      \draw (\xrack-0.66, \y+\dy-0.04) -- (\xrack-0.52, \y+\dy+0.08);
    }
    % Right block
    \draw (\xrack+0.12, \y-0.16) rectangle (\xrack+0.26, \y+0.16);
    \foreach \dy in {-0.12, 0, 0.12} {
      \draw (\xrack+0.12, \y+\dy-0.04) -- (\xrack+0.26, \y+\dy+0.08);
    }
  }

  % Sector gear and slotted arm centered at O1
  \coordinate (ArmEnd) at (\armang:2.35);
  \draw[fill=white]
    (128:\Rgear) arc (128:205:\Rgear)
    .. controls (-0.85, -0.45) and (-0.35, -0.25) .. ($(O1)+(\armang-90:0.16)$)
    -- ($(ArmEnd)+(\armang-90:0.16)$) arc (\armang-90:\armang+90:0.16)
    -- ($(O1)+(\armang+90:0.16)$)
    .. controls (-0.1, 0.35) and (-0.2, 0.95) .. (128:\Rgear);

  % Inner dashed pitch arc on sector gear
  \draw[dashed] (132:{\Rgear-0.14}) arc (132:202:{\Rgear-0.14});

  % Inner slot in arm
  \coordinate (S1) at (\armang:0.85);
  \coordinate (S2) at (\armang:2.18);
  \draw
    ($(S1)+(\armang+90:0.06)$) -- ($(S2)+(\armang+90:0.06)$)
    arc (\armang+90:\armang-90:0.06) -- ($(S1)+(\armang-90:0.06)$)
    arc (\armang-90:\armang-270:0.06) -- cycle;

  % Fixed supports at O1 and O (ceiling-style above)
  \draw (-0.18, 0.34) -- (O1) -- (0.04, 0.30);
  \draw (-0.30, 0.36) -- (0.16, 0.28);
  \foreach \t in {0, 0.2, 0.4, 0.6, 0.8, 1.0} {
    \coordinate (Hpt) at ($(-0.30, 0.36)!\t!(0.16, 0.28)$);
    \draw (Hpt) -- ++(45:0.12);
  }
  \draw[fill=white] (O1) circle (2.4pt);

  \draw (\Dist-0.12, 0.35) -- (O) -- (\Dist+0.12, 0.35);
  \draw (\Dist-0.25, 0.35) -- (\Dist+0.25, 0.35);
  \foreach \x in {-0.20, -0.08, 0.04, 0.16} {
    \draw (\Dist+\x, 0.35) -- (\Dist+\x+0.08, 0.46);
  }

  % Crank OA in front
  \pgfmathsetmacro{\angOA}{180+\cranktheta}
  \draw[fill=white]
    ($(O)+(\angOA+90:0.06)$) -- ($(A)+(\angOA+90:0.06)$)
    arc (\angOA+90:\angOA+270:0.06) -- ($(O)+(\angOA-90:0.06)$) -- cycle;
  \draw[fill=white] (O) circle (2.4pt);
  \fill (A) circle (1.8pt);

  % Horizontal dashed line & angle theta at O
  \draw[dashed] (O) -- (\Dist-0.95, 0);
  \draw ($(O)+(180:0.42)$) arc (180:{180+\cranktheta}:0.42);
  \node at ($(O)+(198:0.60)$) {$\theta$};

  % Rotation arrow n on crank OA
  \draw[->] ($(O)+(255:0.28)$) arc (255:190:0.28);
  \node[below right, inner sep=1pt] at ($(O)+(260:0.30)$) {$n$};

  % Vertical dashed reference lines and dimensions 152cm, 183cm
  \draw[dashed] (\xrack, -1.65) -- (\xrack, -2.05);
  \draw[dashed] (0, 0) -- (0, -2.05);
  \draw[dashed] (\Dist, 0) -- (\Dist, -2.05);

  \draw (\xrack, -2.1) -- (\xrack, -1.75);
  \draw (0, -2.1) -- (0, -1.75);
  \draw (\Dist, -2.1) -- (\Dist, -1.75);
  \draw[<->] (\xrack, -1.92) -- (0, -1.92) node[midway, above, inner sep=3pt] {$152\,\mathrm{cm}$};
  \draw[<->] (0, -1.92) -- (\Dist, -1.92) node[midway, above, inner sep=3pt] {$183\,\mathrm{cm}$};

  % Labels O1, O, A
  \node[above right, inner sep=2pt] at ($(O1)+(0.05, 0.02)$) {$O_1$};
  \node[right, inner sep=3.5pt] at (O) {$O$};
  \node[below, inner sep=3pt] at ($(A)+(0.08,-0.16)$) {$A$};
\end{tikzpicture}

\end{center}
\end{problem}
\clearpage


\begin{problem}{5.1.6}
图 5.7 中机构在竖直平面内运动, 各部件的尺寸如图所示. 某时刻在图示位置, 杆 OA 处于水平位置, 绕 O 点的角速度 $\Omega = 2 \, rad/s$ , 求此时部件 C 的角速度 $\omega$ 及杆 AB 的 B 端的速度的大小.
\begin{center}
\begin{tikzpicture}[>=Stealth, line width=0.8pt]
  % Scale: 1 cm in problem = 0.22 cm in TikZ
  % O=(0, 3.08), A=(1.76, 3.08), O1=(2.20, 0), B=(5.28, 0.44), alpha=-36.87 deg
  \def\s{0.22}
  \coordinate (O) at (0, {14*\s});
  \coordinate (A) at ({8*\s}, {14*\s});
  \coordinate (O1) at ({10*\s}, 0);
  \coordinate (B) at ({24*\s}, {2*\s});
  \def\alphaang{-36.87}

  % Reference dashed lines
  \draw[dashed] (O) -- (0, -0.35);
  \draw[dashed] (-0.35, 0) -- (O1);

  % Dimension 14cm (vertical)
  \draw (-0.85, 0) -- (-0.42, 0);
  \draw (-0.85, {14*\s}) -- (-0.42, {14*\s});
  \draw[<->] (-0.7, 0) -- (-0.7, {14*\s}) node[midway, left, inner sep=2pt] {$14\,\mathrm{cm}$};

  % Dimension 10cm (horizontal)
  \draw (0, -0.45) -- (0, -0.82);
  \draw ({10*\s}, -0.45) -- ({10*\s}, -0.82);
  \draw[<->] (0, -0.68) -- ({10*\s}, -0.68) node[midway, above, inner sep=2pt] {$10\,\mathrm{cm}$};

  % Dimension 8cm (above OA)
  \draw (0, {14*\s+0.25}) -- (0, {14*\s+0.65});
  \draw ({8*\s}, {14*\s+0.25}) -- ({8*\s}, {14*\s+0.65});
  \draw[<->] (0, {14*\s+0.50}) -- ({8*\s}, {14*\s+0.50}) node[midway, above, inner sep=2pt] {$8\,\mathrm{cm}$};

  % Wall support at O
  \draw (-0.22, {14*\s+0.28}) -- (-0.22, {14*\s-0.28});
  \foreach \y in {-0.22, -0.10, 0.02, 0.14, 0.26} {
    \draw (-0.22, {14*\s+\y}) -- (-0.34, {14*\s+\y-0.10});
  }
  \draw (-0.22, {14*\s+0.16}) -- (0, {14*\s+0.16}) arc (90:-90:0.16) -- (-0.22, {14*\s-0.16});

  % Rod AB (dashed inside/before sleeve, solid protruding tip at B)
  \coordinate (Sstart) at ($(A)!0.26!(B)$);
  \coordinate (Send) at ($(A)!0.76!(B)$);
  \draw[dashed] ($(A)+(\alphaang+90:0.06)$) -- ($(Send)+(\alphaang+90:0.06)$);
  \draw[dashed] ($(A)+(\alphaang-90:0.06)$) -- ($(Send)+(\alphaang-90:0.06)$);
  \draw[fill=white]
    ($(Send)+(\alphaang+90:0.05)$) -- ($(B)+(\alphaang+90:0.05)$)
    arc (\alphaang+90:\alphaang-90:0.05) -- ($(Send)+(\alphaang-90:0.05)$) -- cycle;

  % Sleeve on C
  \draw
    ($(Sstart)+(\alphaang+90:0.14)$) -- ($(Send)+(\alphaang+90:0.14)$) --
    ($(Send)+(\alphaang-90:0.14)$) -- ($(Sstart)+(\alphaang-90:0.14)$) -- cycle;

  % Triangular plate C connected to sleeve and pivot O1
  \coordinate (C1) at ($(Sstart)!0.20!(Send)+(\alphaang-90:0.14)$);
  \coordinate (C2) at ($(Sstart)!0.72!(Send)+(\alphaang-90:0.14)$);
  \draw (C1) -- ($(O1)+(143:0.18)$) arc (143:323:0.18) -- (C2);

  % Ground bracket at O1
  \draw ($(O1)+(-0.22,-0.32)$) -- ($(O1)+(-0.22,-0.02)$) arc (180:0:0.22) -- ($(O1)+(0.22,-0.32)$);
  \draw ($(O1)+(-0.38,-0.32)$) -- ($(O1)+(0.38,-0.32)$);
  \foreach \x in {-0.30, -0.18, -0.06, 0.06, 0.18, 0.30} {
    \draw ($(O1)+(\x,-0.32)$) -- ($(O1)+(\x-0.10,-0.44)$);
  }
  \draw[fill=white] (O1) circle (2.2pt);

  % Crank OA
  \draw[fill=white]
    (0, {14*\s+0.08}) -- ({8*\s}, {14*\s+0.08})
    arc (90:-90:0.08) -- (0, {14*\s-0.08}) -- cycle;
  \draw[fill=white] (O) circle (2.4pt);
  \draw[fill=white] (A) circle (2.4pt);

  % Angular velocity arrows Omega and omega
  \draw[->] ($(O)+(0.8, 0.32)$) arc (35:-35:0.55) node[below right, inner sep=1pt] {$\Omega$};
  \draw[->] ($(O1)+(-25:0.32)$) arc (-25:105:0.32) node[above, inner sep=2pt] {$\omega$};

  % Dimension 20cm (along AB)
  \draw ($(A)+(\alphaang+90:0.22)$) -- ($(A)+(\alphaang+90:0.82)$);
  \draw ($(B)+(\alphaang+90:0.22)$) -- ($(B)+(\alphaang+90:0.82)$);
  \draw[<->] ($(A)+(\alphaang+90:0.65)$) -- ($(B)+(\alphaang+90:0.65)$)
    node[midway, above right, inner sep=1pt] {$20\,\mathrm{cm}$};

  % Dimension 10cm (perpendicular distance from O1 to AB)
  \draw[dashed] ($(O1)+(\alphaang:0.45)$) -- ($(O1)+(\alphaang:2.85)$);
  \draw ($(B)+(\alphaang:0.18)$) -- ($(B)+(\alphaang:0.65)$);
  \coordinate (Pdim) at ($(O1)+(\alphaang:2.65)$);
  \coordinate (Bdim) at ($(Pdim)+(\alphaang+90:{10*\s})$);
  \draw[<->] (Pdim) -- (Bdim) node[midway, below right, inner sep=2pt] {$10\,\mathrm{cm}$};

  % Labels
  \node[above left, inner sep=2pt] at ($(O)+(-0.15, 0.2)$) {$O$};
  \node[above right, inner sep=3pt] at (A) {$A$};
  \node[above right, inner sep=3pt] at (B) {$B$};
  \node[above left, inner sep=1pt] at ($(O1)+(-0.22, 0.03)$) {$O_1$};
  \node at ($(O1)+(0.45, 0.72)$) {$C$};
\end{tikzpicture}

\end{center}
\end{problem}
\clearpage


\begin{problem}{5.1.7}
若除上题所给条件以外, 还知道此时 $OA$ 杆绕 $O$ 点的角加速度 $\Omega = 0$ , 求此时部件 $C$ 的角加速度 $\dot{\omega}$ 及 $B$ 点的加速度 $a_B$ .
\end{problem}
\clearpage


\begin{problem}{5.1.8}
某瞬时刚体绕通过坐标原点的某轴转动, 刚体上一点 $M_{1}(1,0,1)$ 的速度 $v_{1} = 4\mathrm{m}\cdot \mathrm{s}^{-1}$ , 它与 $x$ 轴所成的角 $\alpha_{1} = 45^{\circ}$ ; 另一点 $M_{2}(3,4,0)$ 的速度 $v_{2}$ 与 $x$ 轴成 $\alpha_{2}$ 角, 且 $\cos \alpha_{2} = -0.8$ . 试求此刻刚体的角速度 $\omega$ 及 $M_{2}$ 点的速度的大小.
\end{problem}
\clearpage


\begin{problem}{5.1.9}
半径为 $b$ 的圆筒在半径为 $a$ 的固定圆柱上做纯滚动。若相切的直线从静止开始以恒定的角加速度 $\beta$ 围绕固定圆柱的轴顺时针转动，图5.12中画的是 $t = 0$ 时的位置。求圆筒上 $P, C$ 两点 $t$ 时刻的速度和加速度（均按圆筒的切向和法向分解）。
\begin{center}
\begin{tikzpicture}[>=Stealth, line width=0.8pt]
  \def\b{2.15}
  \def\a{0.82}
  \coordinate (C) at (0,0);
  \coordinate (O) at (0, {\b-\a});
  \coordinate (P) at (0, -\b);

  % Outer circle (radius b) and inner circle (radius a)
  \draw (C) circle (\b);
  \draw (O) circle (\a);

  % Vertical diameter through O, C, P
  \draw (0, \b) -- (P);

  % Radius arrow a (from O outward)
  \draw[->] (O) -- ($(O)+(35:\a)$) node[midway, below right, inner sep=1pt] {$a$};

  % Radius arrow b (from outer rim inward to C)
  \draw[->] (205:\b) -- (C) node[midway, below right, inner sep=1pt] {$b$};

  % Angular acceleration arrow beta (clockwise, upper right)
  \draw[->] (42:{\b+0.32}) arc (42:18:{\b+0.32});
  \node[above right, inner sep=1pt] at (32:{\b+0.36}) {$\beta$};

  % Labels O, C, P
  \node[left, inner sep=2.5pt] at (O) {$O$};
  \node[right, inner sep=2.5pt] at (C) {$C$};
  \node[below, inner sep=3pt] at (P) {$P$};
\end{tikzpicture}

\end{center}
\begin{center}
\begin{tikzpicture}[>=Stealth]
  % Fixed cylinder of radius a = 0.85 centered at O, rolling cylinder of radius b = 1.90
  \coordinate (O) at (0,0);
  \coordinate (C) at (-0.860,-0.602);
  \coordinate (T) at (0.696,0.488);
  \coordinate (Q) at (-2.416,-1.692);
  \coordinate (P) at (-1.950,-2.158);
  \coordinate (U) at (0.214,0.937);

  % Coordinate axes
  \draw[->, thick] (0,-3.10) -- (0,2.20) node[left] {$y$};
  \draw[->, thick] (0,0) -- (2.15,0) node[below=2pt] {$x$};
  \node[below right=1pt] at (O) {$O$};

  % Fixed inner cylinder
  \draw[thick] (O) circle (0.85);

  % Rolling cylinder at t = 0 (dashed) and time t (solid)
  \draw[thick, dashed] (0,-1.05) circle (1.90);
  \draw[thick] (C) circle (1.90);

  % Radial line T - O - C - Q and diameter segment P - C - U
  \draw[thick] (T) -- (Q);
  \draw[thick] (P) -- (U);

  % Angle theta from +y to OT
  \draw (0,0.45) arc (90:35:0.45);
  \node at (62:0.65) {$\theta$};

  % Angle (a/b)*theta between CQ and CP
  \draw ($(C)+(215:0.80)$) arc (215:235:0.80);
  \node[below right=1pt] at ($(C)+(236:0.98)$) {$\frac{a}{b}\theta$};

  % Point labels
  \node[above left=1pt] at (C) {$C$};
  \node[below left=1pt] at (Q) {$Q$};
  \node[below left=1pt] at (P) {$P$};
\end{tikzpicture}

\end{center}
\end{problem}
\clearpage


\begin{problem}{5.1.10}
证明刚体上任何两点的速度在它们连线上的投影相等(速度投影定理).
\end{problem}
\clearpage


\begin{problem}{5.1.11}
曲柄 $OA$ 在图 5.14 所示位置 ( $OA$ 与 $DC$ 的夹角为 $60^{\circ}$ . $BD$ 与 $OA$ 垂直) 绕定轴 $O$ 转动的角速度 $\omega = 4\mathrm{s}^{-1}$ , $OA = AB = AD = a = 30\mathrm{cm}$ , $BC = b = 40\mathrm{cm}$ , 滑块 $C, D$ 在同一水平直线上, 并限在水平滑道内滑动, 求此时刻滑块 $C, D$ 的速度.
\begin{center}
\begin{tikzpicture}[>=Stealth]
  % Unit a = 1.45 cm; D=(0,0), G=(0,2a), A=(sqrt(3)/2 a, a/2), B=(sqrt(3) a, a)
  \coordinate (D) at (0,0);
  \coordinate (G) at (0,2.90);
  \coordinate (A) at (1.256,0.725);
  \coordinate (B) at (2.511,1.450);
  \coordinate (O) at (1.981,-0.531);
  \coordinate (C) at (3.790,0);
  \coordinate (M) at (1.674,0);

  % Triangle GBD, crank GO, rod BC
  \draw[thick] (D) -- (G) -- (B) -- (D);
  \draw[thick] (G) -- (O);
  \draw[thick] (B) -- (C);

  % Velocity vectors along DC and dashed connecting segment
  \draw[->, thick] (D) -- (0.95,0) node[midway, below=2pt] {$v_D$};
  \draw[->, thick] (C) -- (2.84,0) node[midway, below=2pt] {$v_C$};
  \draw[thick, dashed] (0.95,0) -- (2.84,0);

  % Velocity vectors at A and B
  \draw[->, thick] (A) -- ($(A)+(30:0.80)$) node[midway, below right=2pt] {$v_A$};
  \draw[->, thick] (B) -- ($(B)+(60:0.95)$) node[above=1pt] {$v_B$};

  % Dashed extensions at B
  \draw[thick, dashed] (B) -- ($(B)+(30:0.75)$);
  \draw[thick, dashed] (B) -- ($(B)+(131.4:0.65)$);

  % Right-angle marks at D, A, and B
  \draw (0.20,0) -- (0.20,0.20) -- (0,0.20);
  \draw ($(A)+(30:0.18)$) -- ($(A)+(30:0.18)+(-60:0.18)$) -- ($(A)+(-60:0.18)$);
  \draw ($(B)+(60:0.18)$) -- ($(B)+(60:0.18)+(150:0.18)$) -- ($(B)+(150:0.18)$);

  % Angle 60 deg at M and 30 deg at B
  \draw ($(M)+(180:0.42)$) arc (180:120:0.42);
  \node at ($(M)+(150:0.68)$) {$60^\circ$};
  \draw ($(B)+(30:0.46)$) arc (30:60:0.46);
  \node[right=1pt] at ($(B)+(44:0.66)$) {$30^\circ$};

  % Point and length labels
  \node[above=2pt] at (G) {$G$};
  \node[left=2pt] at (D) {$D$};
  \node[below right=1pt] at (O) {$O$};
  \node[right=2pt] at (C) {$C$};
  \node[above left=2pt] at (A) {$A$};
  \node[below=3pt] at (B) {$B$};
  \node[right=2pt] at ($(G)!0.64!(A)$) {$\sqrt{3}a$};
  \node[above right=2pt] at ($(G)!0.42!(B)$) {$2a$};
\end{tikzpicture}

\end{center}
\begin{center}
\begin{tikzpicture}[>=Stealth]
  % Unit a = 1.40 cm
  \coordinate (D) at (0,0);
  \coordinate (A) at (1.212,0.700);
  \coordinate (B) at (2.425,1.400);
  \coordinate (O) at (1.912,-0.512);
  \coordinate (C) at (3.660,0);
  \coordinate (M) at (1.617,0);

  % Horizontal dashed axis
  \draw[thick, dashed] (-0.65,0) -- (4.30,0);

  % Slider D and guide rails
  \fill[white] (-0.32,-0.16) rectangle (0.32,0.16);
  \draw[thick] (-0.32,-0.16) rectangle (0.32,0.16);
  \draw[thick] (-0.55,0.20) -- (0.55,0.20);
  \draw[thick] (-0.55,-0.20) -- (0.55,-0.20);
  \foreach \x in {-0.45,-0.25,-0.05,0.15,0.35} {
    \draw (\x,0.20) -- (\x+0.14,0.34);
    \draw (\x,-0.34) -- (\x+0.14,-0.20);
  }
  \node[above=12pt] at (D) {$D$};

  % Slider C and guide rails
  \fill[white] (3.34,-0.16) rectangle (3.98,0.16);
  \draw[thick] (3.34,-0.16) rectangle (3.98,0.16);
  \draw[thick] (3.11,0.20) -- (4.21,0.20);
  \draw[thick] (3.11,-0.20) -- (4.21,-0.20);
  \foreach \x in {3.21,3.41,3.61,3.81,4.01} {
    \draw (\x,0.20) -- (\x+0.14,0.34);
    \draw (\x,-0.34) -- (\x+0.14,-0.20);
  }
  \node[above=12pt] at (C) {$C$};

  % Fixed pivot support at O
  \draw[thick, fill=white] ($(O)+(-0.24,-0.28)$) -- ($(O)+(-0.18,0)$)
    arc (180:0:0.18) -- ($(O)+(0.24,-0.28)$) -- cycle;
  \draw[thick] ($(O)+(-0.36,-0.28)$) -- ($(O)+(0.36,-0.28)$);
  \foreach \x in {-0.28,-0.14,0,0.14,0.28} {
    \draw ($(O)+(\x-0.10,-0.40)$) -- ($(O)+(\x,-0.28)$);
  }
  \node[left=6pt] at (O) {$O$};

  % Links OA, DB, BC
  \draw[thick] (O) -- (A);
  \draw[thick] (D) -- (B);
  \draw[thick] (B) -- (C);

  % Right-angle mark at A
  \draw ($(A)+(30:0.18)$) -- ($(A)+(30:0.18)+(-60:0.18)$) -- ($(A)+(-60:0.18)$);

  % Pin joints at D, A, B, O, C
  \foreach \pt in {D,A,B,O,C} {
    \fill[white] (\pt) circle (0.08);
    \draw[thick] (\pt) circle (0.08);
    \fill (\pt) circle (0.02);
  }

  % Angle 60 deg
  \draw ($(M)+(180:0.42)$) arc (180:120:0.42);
  \node at ($(M)+(150:0.65)$) {$60^\circ$};

  % Rotation arrow omega around O
  \draw[->, thick] ($(O)+(150:0.46)$) arc (150:25:0.46);
  \node[right=1pt] at ($(O)+(20:0.48)$) {$\omega$};

  % Labels
  \node[above left=2pt] at (A) {$A$};
  \node[above right=2pt] at (B) {$B$};
  \node[above left=2pt] at ($(D)!0.55!(A)$) {$a$};
  \node[above left=2pt] at ($(A)!0.5!(B)$) {$a$};
  \node[right=2pt] at ($(A)!0.38!(O)$) {$a$};
  \node[right=2pt] at ($(B)!0.44!(C)$) {$b$};
\end{tikzpicture}

\end{center}
\end{problem}
\clearpage


\begin{problem}{5.1.12}
一个半径为 $r$ 的圆柱沿一平面做无滑滚动，其轴线 $O$ 保持速度 $v_{0}$ 不变，一块长度为 $l$ 的矩形平板其 $B$ 边以铰链接在圆柱的边缘上，另一端 $A$ 边搁在上述平面上随圆柱运动。若 $t = 0$ 时 $\varphi = 0$ ，求 $A$ 边的速度与时间的关系，并取适当的坐标，取最方便的基点，写出 $AB$ 板的运动学方程。
\end{problem}
\clearpage


\begin{problem}{5.1.13}
求上题中 $AB$ 板的空间极迹和本体极迹方程(可用 $t$ 作为参数).
\end{problem}
\clearpage


\begin{problem}{5.1.14}
证明均质的质量为 $m$ 的直棒的动能 $T = \frac{1}{6} m(u \cdot u + u \cdot v + v \cdot v)$ , 其中 $u, v$ 分别是棒两端的速度.
\end{problem}
\clearpage


\begin{problem}{5.1.15}
如图 5.17 曲柄 OA 以匀角速 $n_{0}=50r/min$ 绕定轴 O 转动，定齿轮的齿数为 $Z_{4}$ ，齿数 $Z_{1}, Z_{2}$ 的两齿轮紧固成一体装在曲柄 OA 上的 B 点。求装在 A 端的齿数 $Z_{3}$ 的齿轮转速。
\begin{center}
\begin{tikzpicture}[>=Stealth]
  \coordinate (O) at (0,0);
  \coordinate (B) at (1.78,0);
  \coordinate (A) at (3.78,0);

  % Outer middle gear Z2 (drawn first so Z4 cleanly occludes overlap)
  \draw[thick] (B) circle (1.00);

  % Fixed left gear Z4
  \fill[white] (O) circle (1.00);
  \draw[thick] (O) circle (1.00);

  % Inner middle gear Z1 and right gear Z3
  \draw[thick] (B) circle (0.78);
  \draw[thick] (A) circle (1.00);

  % Ground hatching under gear Z4
  \draw[thick] (-0.42,-1.00) -- (0.42,-1.00);
  \foreach \x in {-0.32,-0.16,0,0.16,0.32} {
    \draw (\x-0.12,-1.18) -- (\x,-1.00);
  }
  \draw[thick] (0.42,-1.00) -- (0.50,-1.18);

  % Crank OA (rounded capsule)
  \filldraw[fill=white, draw=black, thick]
    (0,0.13) -- (3.78,0.13) arc (90:-90:0.13) -- (0,-0.13) arc (270:90:0.13) -- cycle;

  % Center dots and labels O, B, A
  \foreach \pt/\lbl in {O/O, B/B, A/A} {
    \fill (\pt) circle (0.035);
    \node[below=5pt] at (\pt) {$\lbl$};
  }

  % Gear labels
  \node at (-0.40,0.58) {$Z_4$};
  \node[above left=1pt] at (1.15,0.92) {$Z_2$};
  \node at (1.44,0.50) {$Z_1$};
  \node at (4.12,0.58) {$Z_3$};
\end{tikzpicture}

\end{center}
\end{problem}
\clearpage


\begin{problem}{5.1.16}
半径为 $b$ 的小圆柱在一个较大的半径为 $a$ 的圆柱内部纯滚动，如外面的圆柱围绕它的轴(固定轴)以角速度 $\Omega$ 转动，两圆柱的轴所在的平面以角速度 $\omega$ 绕固定轴转动，设 $\Omega, \omega$ 都是正的，求小圆柱的角速度。
\end{problem}
\clearpage


\begin{problem}{5.1.17}
若做定点转动的刚体的运动学方程为

\[
\varphi = a t, \quad \psi = b t, \quad \theta = c
\]

式中 a、b、c 均为常量， $\varphi$ 、 $\psi$ 、 $\theta$ 为欧拉角、求角速度、角加速度在动、静坐标系中的分量。
\end{problem}
\clearpage


\begin{problem}{5.1.18}
刚体做上题的运动, 求刚体上 $(1,0,0)$ 点及位于空间 $(1,0,0)$ 的刚体上的点在 $t$ 时刻的速度和加速度.
\end{problem}
\clearpage


\begin{problem}{5.1.19}
刚体做 5.1.17 题所给的定点转动, 求本体极面和空间极面满足的方程.
\end{problem}
\clearpage


\begin{problem}{5.1.20}
开始固连于刚体的动坐标系 $Oxyz$ 与静坐标系 $O\xi \eta \zeta$ 完全重合 $(x, y, z$ 轴分别与 $\xi, \eta, \zeta$ 轴重合). 动系先绕 $\xi$ 轴转 $\theta$ 角, 再绕新的 $y$ 轴转 $\psi$ 角, 最后绕新的 $z$ 轴转 $\varphi$ 角. 刚体的位置用此 $\theta, \psi, \varphi$ 来描述. 像推导欧拉运动学方程那样, 分别写出做定点转动的刚体的角速度在固连于刚体的 $xyz$ 坐标系的三个分量, 与在 $\xi \eta \zeta$ 静坐标系的三个分量, 与这里定义的 $\theta, \psi, \varphi$ 三个角及其对时间的导数的关系.
\end{problem}
\clearpage


\section{刚体的相对运动}

\begin{problem}{5.2.1}
一个半径为 $a$ 、圆心为 $A$ 的圆柱在一个水平面上做无滑滚动，圆心 $A$ 以恒定速度 $u$ 运动；另一半径为 $b$ 、圆心为 $B$ 的第二个圆柱（它的轴与第一个圆柱的轴平行）在第一个圆柱上做无滑滚动， $B$ 点相对于固连于第一个圆柱的坐标系以恒定的速率 $v$ 运动。 $P$ 点是第二个圆柱边缘上的点，若 $t = 0$ 时， $A, B, P$ 在同一铅垂线上。求 $P$ 点的速度、加速度和时间的关系。
\end{problem}
\clearpage


\begin{problem}{5.2.2}
A、B、C 三个半径均为 r 的圆盘，紧挨着安装在一曲柄上．除 A 盘不动外，B、C 两盘可分别绕 $O_{2}$ 、 $O_{3}$ 转动，曲柄围绕 $O_{1}$ 转动．整个机构在同一平面上．B、C 盘均做滚动．若 B 盘的转动角速度为 $\omega(t)$ ．求：

(1) C 盘的转动角速度；

(2) t 时刻位于 B、C 两盘相接触的两点 P、Q 的加速度的大小；

(3) 若 t=0 时, P 点在图示位置, 求任何时刻 P 点的加速度的大小.
\begin{center}
\begin{tikzpicture}[>=Stealth]
  \coordinate (O1) at (0,0);
  \coordinate (O2) at (1.497,1.170);
  \coordinate (O3) at (2.994,2.339);
  \coordinate (Mbc) at ($(O2)!0.5!(O3)$);

  % Disks A, B, C
  \draw[thick] (O1) circle (0.95);
  \draw[thick] (O2) circle (0.95);
  \draw[thick] (O3) circle (0.95);

  % Radius lines labeled r
  \foreach \c in {O1,O2,O3} {
    \draw[thick] (\c) -- ($(\c)+(98:0.95)$) node[midway, right=1pt] {$r$};
  }

  % Ground base under disk A
  \draw[thick] (-0.45,-0.95) -- (0.45,-0.95);
  \foreach \x in {-0.34,-0.17,0,0.17,0.34} {
    \draw (\x-0.14,-1.18) -- (\x,-0.95);
  }
  \draw[thick] (0.45,-0.95) -- (0.56,-1.18);

  % Crank O1-O2-O3
  \draw[thick]
    ($(O1)+(128:0.13)$) -- ($(O3)+(128:0.13)$)
    arc (128:-52:0.13) -- ($(O1)+(-52:0.13)$)
    arc (-52:-232:0.13) -- cycle;

  % Center dots and labels O1, O2, O3
  \foreach \c/\lbl in {O1/O_1, O2/O_2, O3/O_3} {
    \fill (\c) circle (0.035);
    \node[below right=4pt] at (\c) {$\lbl$};
  }

  % Disk labels A, B, C
  \node at (-1.12,-0.25) {$A$};
  \node at ($(O2)+(145:1.12)$) {$B$};
  \node at ($(O3)+(145:1.12)$) {$C$};

  % Contact point labels P and Q
  \node at ($(Mbc)+(-142:0.22)+(-52:0.29)$) {$P$};
  \node at ($(Mbc)+(38:0.22)+(128:0.29)$) {$Q$};

  % Rotation arrow omega(t) for disk B
  \draw[->, thick] ($(O2)+(-25:1.15)$) arc (-25:-75:1.15);
  \node[right=2pt] at ($(O2)+(-48:1.15)$) {$\omega(t)$};
\end{tikzpicture}

\end{center}
\end{problem}
\clearpage


\begin{problem}{5.2.3}
一半径为 $a$ 的圆盘 $A$ 以等角速 $\omega_{A}$ 绕其固定中心 $O$ 转动，另一半径为 $b$ 的圆

盘 B 以等角速 $\omega_{B}$ 绕其中心 $O'$ 在前一圆盘外侧纯滚动(如图 5.20 所示). 求:

(1) 两盘的连杆 $OO'$ 的转动角速度 $\Omega$ ;

(2) B 盘在 A 盘上转动回到 A 盘上同一接触点所需的最短时间 $t_{1}$ ;

(3) B 盘上与 A 盘相接触的点又回到与 A 盘相接触的位置所需的最短时间 $t_{2}$ ;

(4) $O'$ 点回到空间原来位置所需的最短时间 $t_3$ ;

（5）如 $a = 5\mathrm{cm}, b = 3\mathrm{cm}, \omega_A = 20\pi \mathrm{s}^{-1}, \omega_B = 36\pi \mathrm{s}^{-1}, B$ 盘所有点回到空间原来位置所需的最小时间 $t_4, A, B$ 盘所有点回到空间原来位置所需的最小时间 $t_5, A$ 盘所有点及 $O'$ 点回到空间原来位置所需的最小时间 $t_6$ .
\begin{center}
\begin{tikzpicture}[>=Stealth]
  \coordinate (O) at (0,0);
  \coordinate (Op) at (2.144,1.558);

  % Disk A and Disk B
  \draw[thick] (O) circle (1.75);
  \draw[thick] (Op) circle (0.90);

  % Radius lines a and b
  \draw[thick] (O) -- (200:1.75) node[midway, above=1pt] {$a$};
  \draw[thick] (Op) -- ($(Op)+(0,-0.90)$) node[midway, right=1pt] {$b$};

  % Fixed pivot support at O
  \draw[thick, fill=white] ($(O)+(-0.28,-0.34)$) -- ($(O)+(-0.20,0.02)$)
    arc (180:0:0.20) -- ($(O)+(0.28,-0.34)$) -- cycle;
  \draw[thick] ($(O)+(-0.46,-0.34)$) -- ($(O)+(0.46,-0.34)$);
  \foreach \x in {-0.34,-0.17,0,0.17,0.34} {
    \draw (\x-0.12,-0.50) -- (\x,-0.34);
  }

  % Connecting rod OO'
  \filldraw[fill=white, draw=black, thick]
    ($(O)+(126:0.14)$) -- ($(Op)+(126:0.14)$)
    arc (126:-54:0.14) -- ($(O)+(-54:0.14)$)
    arc (-54:-234:0.14) -- cycle;

  % Center dots and labels O, O'
  \fill (O) circle (0.035);
  \fill (Op) circle (0.035);
  \node[above left=2pt] at ($(O)+(-0.14,0.14)$) {$O$};
  \node[above right=2pt] at ($(Op)+(0.14,0.06)$) {$O'$};

  % Disk labels A, B
  \node at (-52:1.95) {$A$};
  \node at ($(Op)+(-22:1.08)$) {$B$};

  % Angular velocity arrows omega_A and omega_B
  \draw[->, thick] (115:1.30) arc (115:155:1.30);
  \node[right=2pt] at (130:1.28) {$\omega_A$};

  \draw[->, thick] ($(Op)+(105:1.08)$) arc (105:150:1.08);
  \node[below=2pt] at ($(Op)+(122:0.94)$) {$\omega_B$};
\end{tikzpicture}

\end{center}
\end{problem}
\clearpage


\begin{problem}{5.2.4}
高为 $h$ 、顶角为 $2\alpha$ 的正圆锥在一水平面上绕顶点作纯滚动，若已知其几何对称轴以恒定的角速度 $\Omega$ 绕竖直轴转动．求此刻圆锥底面上最高点 $A$ 的速度和加速度．
\end{problem}
\clearpage


\begin{problem}{5.2.5}
直杆 A 铰接于 U 形管夹板的 $OO'$ 轴上，该板连在旋转的竖直轴的端点，竖直轴以图 5.23 所示的角速度 $\Omega$ 转动，杆 A 与竖直轴 z 的夹角 $\theta$ 以不变的速率 $\dot{\theta} = -p$ 变化。求杆 A 的角速度 $\omega$ 与角加速度 $\beta$ 。
\begin{center}
\begin{tikzpicture}[>=Stealth,
  x={(-0.56cm,-0.45cm)}, y={(0.86cm,-0.25cm)}, z={(0cm,1.0cm)}]

  % Origin O'' and coordinate axes x, y, z
  \coordinate (Opp) at (0,0,0);
  \draw[->, thick] (Opp) -- (2.2,0,0) node[below left=1pt] {$x$};
  \draw[->, thick] (Opp) -- (0,2.2,0) node[right=1pt] {$y$};
  \draw[thick] (Opp) -- (0,0,2.70);
  \node[below=2pt] at (Opp) {$O''$};

  % Rod A tilting in xz-plane (angle theta = 27 deg from -z toward +x)
  \coordinate (H) at (0,0,2.75);
  \coordinate (T) at (1.203,0,0.389);
  \coordinate (T_ext) at (1.42,0,-0.04);

  % Angle theta arc between -z axis and rod A
  \draw[thick] (0,0,1.45) .. controls (0.28,0,1.45) and (0.48,0,1.50) .. (0.64,0,1.58);
  \node at (0.34,0,1.28) {$\theta$};

  % Rod A (3D cylinder with dashed centerline)
  \filldraw[fill=white, draw=black, thick]
    ($(H)+(-0.10cm,0.045cm)$) -- ($(T)+(-0.10cm,0.045cm)$)
    .. controls ($(T)+(-0.08cm,-0.08cm)$) and ($(T)+(0.08cm,-0.11cm)$) .. ($(T)+(0.10cm,-0.045cm)$)
    -- ($(H)+(0.10cm,-0.045cm)$) -- cycle;
  \draw[thick, dashed] (H) -- (T_ext);
  \node[left=4pt] at ($(H)!0.55!(T)$) {$A$};

  % Hinge axis oo' (dashed, parallel to y)
  \draw[thick, dashed] (0,-2.1,2.75) node[left=1pt] {$o$} -- (0,-0.52,2.75);
  \draw[thick, dashed] (0,-0.24,2.75) -- (0,0.24,2.75);
  \draw[thick, dashed] (0,0.52,2.75) -- (0,2.0,2.75) node[below right=1pt] {$o'$};

  % U-shaped bracket (clevis) around H = (0,0,2.75)
  \filldraw[fill=white, draw=black, thick]
    (0.28,-0.52,2.85) -- (0.28,-0.52,2.55) -- (0.28,-0.22,2.55) -- (0.28,-0.22,2.85) -- cycle;
  \filldraw[fill=white, draw=black, thick]
    (0.28,0.22,2.85) -- (0.28,0.22,2.55) -- (0.28,0.52,2.55) -- (0.28,0.52,2.85) -- cycle;
  \filldraw[fill=white, draw=black, thick]
    (0.28,0.52,2.85) -- (0.28,0.52,2.55) -- (-0.28,0.52,2.55) -- (-0.28,0.52,3.02)
    -- (0.28,0.52,3.02) -- cycle;
  \filldraw[fill=white, draw=black, thick]
    (0.28,-0.52,2.85) -- (0.28,0.52,2.85) -- (0.28,0.52,3.02) -- (0.28,-0.52,3.02) -- cycle;
  \filldraw[fill=white, draw=black, thick]
    (0.28,-0.52,3.02) -- (0.28,0.52,3.02) -- (-0.28,0.52,3.02) -- (-0.28,-0.52,3.02) -- cycle;

  % Square vertical stem on top of bracket
  \filldraw[fill=white, draw=black, thick]
    (0.13,-0.13,3.02) -- (0.13,0.13,3.02) -- (0.13,0.13,3.55) -- (0.13,-0.13,3.55) -- cycle;
  \filldraw[fill=white, draw=black, thick]
    (0.13,0.13,3.02) -- (-0.13,0.13,3.02) -- (-0.13,0.13,3.55) -- (0.13,0.13,3.55) -- cycle;
  \filldraw[fill=white, draw=black, thick]
    (0.13,-0.13,3.55) -- (0.13,0.13,3.55) -- (-0.13,0.13,3.55) -- (-0.13,-0.13,3.55) -- cycle;

  % z-axis arrow continuing above square stem
  \draw[->, thick] (0,0,3.55) -- (0,0,4.10) node[right=1pt] {$z$};

  % Rotation arrow Omega around vertical stem
  \draw[->, thick] (-0.30,-0.32,3.26) .. controls (-0.12,-0.45,3.15) and (0.42,-0.15,3.12) .. (0.36,0.34,3.34);
  \node[right=3pt] at (0.36,0.34,3.36) {$\Omega$};
\end{tikzpicture}

\end{center}
\end{problem}
\clearpage


\begin{problem}{5.2.6}
飞机重心 G 做匀速直线运动, 坐标系 xyz 的原点置于 G 点, 飞机相对于此坐标系的角速度为: 滚转 $\omega_{x}=0.5s^{-1}$ , 俯仰 $\omega_{y}=0.3s^{-1}$ , 偏航 $\omega_{z}=0.1s^{-1}$ , 现有一旅客在重心 G 之后 30m, 高于 G 点 2m, 以相对于机身的速度 $v_{r}=1m\cdot s^{-1}$ 沿与 x 轴相反的方向走向机尾, 求此旅客的加速度.
\end{problem}
\clearpage


\begin{problem}{5.2.7}
圆盘以匀角速 $\Omega$ 绕其中心轴 $z$ 转动，架子 $A$ 以匀角速度 $\omega_{2}$ 绕其轴线 $y'$ ( $y' \parallel y$ ) 转动。同时，整个装置以角速度 $\omega_{1}$ 绕 $X$ 轴转动。求架子 $A$ 使圆盘处于图5.25所示位置时圆盘的角速度 $\omega$ 、角加速度 $\beta$ 以及圆盘上最高点 $P$ 的速度和加速度，图中 $xyz$ 坐标轴固连于架子 $A$ 。在此时， $x$ 轴与固定坐标 $X$ 轴平行。设圆盘半径为 $r, y, y'$ 轴间距为 $b, x, X$ 轴间距为 $c, z$ 与 $Z$ 轴间的水平间距可以忽略。
\begin{center}
\begin{tikzpicture}[>=Stealth]
  \coordinate (O) at (0,0);
  \coordinate (P) at (0,1.02);
  \coordinate (Lp) at (0,-1.45);
  \coordinate (ArmL) at (0.94,-0.85);
  \coordinate (ArmR) at (2.92,-1.38);
  \coordinate (OXYZ) at (2.92,-2.52);

  % 1. Right vertical shaft, flange collar, and XYZ axes
  \filldraw[fill=white, draw=black, thick] (2.82,-2.52) rectangle (3.02,-1.38);
  \filldraw[fill=white, draw=black, thick] (2.92,-2.05) ellipse (0.44 and 0.14);
  \filldraw[fill=white, draw=black, thick] (2.92,-2.00) ellipse (0.44 and 0.14);
  \fill[white] (2.82,-2.00) rectangle (3.02,-1.45);
  \draw[thick] (2.82,-2.00) -- (2.82,-1.45);
  \draw[thick] (3.02,-2.00) -- (3.02,-1.45);
  \draw[->, thick] (2.50,-2.16) -- (2.50,-1.84);
  \draw[->, thick] (3.34,-2.16) -- (3.34,-1.84);
  \draw[->, thick] (2.70,-1.72) arc (160:380:0.22 and 0.08);
  \node[right=2pt] at (3.14,-1.74) {$\omega_1$};

  \draw[->, thick] (OXYZ) -- ($(OXYZ)+(-0.98,-0.63)$) node[below left=1pt] {$Y$};
  \draw[->, thick] (OXYZ) -- ($(OXYZ)+(1.15,-0.31)$) node[right=1pt] {$Z$};

  % 2. Horizontal arm (rounded 3D slab from ArmL to ArmR)
  \filldraw[fill=white, draw=black, thick]
    ($(ArmL)+(-0.12,0.34)$) -- ($(ArmR)+(-0.08,0.34)$)
    .. controls ($(ArmR)+(0.32,0.34)$) and ($(ArmR)+(0.44,-0.02)$) .. ($(ArmR)+(0.26,-0.22)$)
    -- ($(ArmL)+(0,-0.28)$) -- cycle;
  \filldraw[fill=white, draw=black, thick]
    ($(ArmL)+(0,0.28)$) -- ($(ArmR)+(0,0.28)$)
    arc (90:-90:0.28) -- ($(ArmL)+(0,-0.28)$) arc (270:90:0.28) -- cycle;

  % X-axis rising from right pivot on ArmR
  \fill[white] (ArmR) circle (0.075);
  \draw[thick] (ArmR) circle (0.075);
  \draw[->, thick] ($(ArmR)+(0,0.075)$) -- (2.92,1.25) node[left=1pt] {$X$};

  % 3. Disk centered at O=(0,0) behind frame A
  \filldraw[fill=white, draw=black, thick]
    (-0.20,1.02) -- (0,1.02) arc (90:270:0.78 and 1.02) -- (-0.20,-1.02)
    arc (270:90:0.78 and 1.02) -- cycle;
  \filldraw[fill=white, draw=black, thick] (O) ellipse (0.78 and 1.02);
  \draw[->, thick] (-0.42,1.16) arc (122:162:1.25);
  \node[above left=1pt] at (-0.64,1.06) {$\Omega$};

  % 4. Shaft segment from ArmL forward-left toward Lp
  \draw[thick] ($(ArmL)+(-0.20,-0.04)$) -- (0.28,-1.21);
  \draw[thick] ($(ArmL)+(-0.10,-0.16)$) -- (0.28,-1.35);

  % 5. Frame A (vertical rounded-end bracket between Lp and O)
  \filldraw[fill=white, draw=black, thick]
    (0.36,0) arc (0:180:0.26 and 0.22) -- (-0.16,-1.45)
    arc (180:360:0.26 and 0.22) -- cycle;
  \filldraw[fill=white, draw=black, thick]
    (0.26,0) arc (0:180:0.26 and 0.22) -- (-0.26,-1.45)
    arc (180:360:0.26 and 0.22) -- cycle;
  \node at (0,-0.28) {$O$};
  \node at (0,-0.84) {$A$};

  % Dashed continuation of y' shaft inside frame A
  \draw[thick, dashed] (0.26,-1.28) -- (Lp);

  % Lower pivot pin at Lp and upper pivot pin at O
  \fill[white] (Lp) circle (0.075);
  \draw[thick] (Lp) circle (0.075);
  \fill (Lp) circle (0.025);
  \fill[white] (O) circle (0.075);
  \draw[thick] (O) circle (0.075);

  % y' axis from Lp forward-left, with omega_2 arrow
  \draw[->, thick] (Lp) -- (-1.28,-2.27) node[below=1pt] {$y'$};
  \draw[->, thick] (-0.98,-1.90) arc (125:345:0.16 and 0.22);
  \node[right=2pt] at (-0.72,-2.10) {$\omega_2$};

  % 6. Axes x, y, z at O
  \draw[->, thick] (0,0.075) -- (0,1.75) node[left=1pt] {$x$};
  \node[above left=2pt] at (-0.04,1.02) {$P$};
  \draw[->, thick] (O) -- (-1.28,-0.82) node[left=1pt] {$y$};
  \draw[->, thick] (0.075,-0.02) -- (1.05,-0.28) node[right=1pt] {$z$};

  % 7. Dimension arrows b and c
  \draw[<->, thick] (-1.05,-0.67) -- (-1.05,-2.12) node[midway, left=1pt] {$b$};
  \draw[<->, thick] (P) -- (2.92,0.24) node[midway, above=1pt] {$c$};
\end{tikzpicture}

\end{center}
\end{problem}
\clearpage


\begin{problem}{5.2.8}
如图 5.26 所示, 圆锥滚子在水平的圆锥环形支座上做纯滚动, 滚子底面半径 $R = 10 \sqrt{2} \, cm$ , 顶角 $2\alpha = 90^{\circ}$ , 滚子中心 A 沿其轨迹运动的速率 $v_{A} = 20 \, cm \cdot s^{-1}$ , 求圆锥滚子上点 B 和 C 的速度和加速度的大小.
\begin{center}
\begin{tikzpicture}[>=Stealth]
  \coordinate (O) at (0,0);
  \coordinate (B) at (1.60,1.60);
  \coordinate (C) at (1.60,-1.60);

  % Conical seat (ground base below)
  \draw[thick] (-2.45,-2.00) -- (-2.00,-2.00) -- (-0.95,-0.95) -- (0.95,-0.95)
    -- (2.00,-2.00) -- (2.45,-2.00);
  \foreach \x in {-0.80,-0.60,-0.40,-0.20,0,0.20,0.40,0.60,0.80} {
    \draw (\x-0.14,-1.13) -- (\x,-0.95);
  }
  \foreach \x in {-2.35,-2.15,-1.95} {
    \draw (\x-0.12,-2.16) -- (\x,-2.00);
  }
  \foreach \x in {2.10,2.30,2.50} {
    \draw (\x-0.12,-2.16) -- (\x,-2.00);
  }

  % Left and right hatched conical roller rims
  \filldraw[pattern=north east lines, draw=black, thick]
    (-1.60,-1.60) -- (-1.60,1.60) -- (-1.28,1.28) -- (-1.28,-1.28) -- cycle;
  \filldraw[pattern=north east lines, draw=black, thick]
    (1.28,-1.28) -- (1.28,1.28) -- (1.60,1.60) -- (1.60,-1.60) -- cycle;

  % Central block and horizontal shaft
  \draw[thick, fill=white] (-0.52,-0.48) rectangle (0.52,0.48);
  \draw[thick] (-1.28,-0.10) rectangle (-0.52,0.10);
  \draw[thick] (0.52,-0.10) rectangle (1.28,0.10);
  \draw[thick] (-1.98,-0.16) rectangle (-1.60,0.16);
  \draw[thick] (1.60,-0.16) rectangle (1.98,0.16);

  % 45-degree dashed generators through O
  \draw[thick, dashed] (-1.28,-1.28) -- (1.28,1.28);
  \draw[thick, dashed] (-1.28,1.28) -- (1.28,-1.28);

  % Angle 2*alpha arc on the left (cleanly placed in gap [-1.28, -0.52])
  \draw[<->, thick] (135:0.74) arc (135:225:0.74);
  \node at (-0.92,-0.42) {$2\alpha$};

  % Labels O, A, B, C
  \node[fill=white, inner sep=0.8pt] at (0,0.24) {$O$};
  \node[right=2pt] at (1.98,0) {$A$};
  \node[right=2pt] at (B) {$B$};
  \node[right=2pt] at (C) {$C$};
\end{tikzpicture}

\end{center}
\begin{center}
\begin{tikzpicture}[>=Stealth]
  \coordinate (O) at (0,0);
  \coordinate (A) at (1.85,0);
  \coordinate (B) at (1.85,1.85);
  \coordinate (C) at (1.85,-1.85);
  \coordinate (D) at (-1.85,-1.85);

  % Dashed cone geometry lines and solid segment BC
  \draw[thick, dashed] (O) -- (A);
  \draw[thick, dashed] (D) -- (B);
  \draw[thick, dashed] (O) -- (C);
  \draw[thick, dashed] (D) -- (C);
  \draw[thick] (B) -- (C);

  % Angular velocity vectors omega_1, omega_2, omega
  \draw[->, thick] (O) -- (0,1.35) node[right=2pt] {$\boldsymbol{\omega}_1$};
  \draw[->, thick] (O) -- (-1.35,0) node[below=2pt] {$\boldsymbol{\omega}_2$};
  \draw[->, thick] (O) -- (-1.35,1.35) node[above=1pt] {$\boldsymbol{\omega}$};
  \draw[thick, dashed] (0,1.35) -- (-1.35,1.35) -- (-1.35,0);

  % Point labels
  \node[below=3pt] at (O) {$O$};
  \node[right=2pt] at (A) {$A$};
  \node[above right=1pt] at (B) {$B$};
  \node[below right=1pt] at (C) {$C$};
\end{tikzpicture}

\end{center}
\end{problem}
\clearpage


\begin{problem}{5.2.9}
图 5.28 是一个双重差动机构, 曲柄 3 绕固定轴 AB 转动、在曲柄上活动地套一行星齿轮 4, 行星齿轮由半径分别为 $r_{1}=5cm$ 和 $r_{2}=2cm$ 的两个锥齿轮固连而成, 两锥齿轮又分别与半径分别为 $R_{1}=10cm$ 、 $R_{2}=5cm$ 的另两个锥齿轮 1、2 啮合, 齿轮 1、2 均绕 AB 轴转动, 均不与曲柄相连, 它们的角速度分别为 $\omega_{1}=4.5rad/s$ , $\omega_{2}=9rad/s$ , 转动方向相同. 求曲柄 3 的角速度 $\omega_{3}$ 和行星齿轮相对于曲柄的角速度 $\omega_{43}$ .
\begin{center}
\begin{tikzpicture}[>=Stealth, thick]
  % Central vertical axis AB
  \draw (0,-1.95) -- (0,1.95);
  \node[right] at (0,1.90) {$A$};
  \node[right] at (0,-1.90) {$B$};
  \node[left] at (-0.06,0) {$O$};

  % Crank 3
  \draw (0,0) -- (2.8,0);
  \node[below] at (0.65,-0.04) {$3$};

  % Upper bevel gear 2 (radius R2)
  \draw (-0.10,0.70) -- (-0.10,1.20) (0.10,0.70) -- (0.10,1.20);
  \draw (-0.10,0.95) -- (-1.05,0.95);
  \node[below] at (-0.60,0.93) {$2$};
  \draw (-1.05,0.95) ++(-42:0.25) -- ++(138:0.50);

  \draw (0.10,0.95) -- (1.02,0.95);
  \node[above] at (0.56,0.97) {$R_2$};
  \draw (1.02,0.95) ++(-138:0.26) -- ++(42:0.52);

  % Lower bevel gear 1 (radius R1)
  \draw (-0.10,-1.50) -- (-0.10,-1.00) (0.10,-1.50) -- (0.10,-1.00);
  \draw (-0.10,-1.25) -- (-2.15,-1.25);
  \node[above] at (-1.15,-1.23) {$1$};
  \draw (-2.15,-1.25) ++(-138:0.25) -- ++(42:0.50);

  \draw (0.10,-1.25) -- (2.12,-1.25);
  \node[below] at (1.12,-1.27) {$R_1$};
  \draw (2.12,-1.25) ++(138:0.26) -- ++(-42:0.52);

  % Planetary gear 4 (r2 and r1 mounted on crank 3)
  % Upper half
  \draw (1.10,0.87) -- (1.10,0.12) -- (2.20,0.12) -- (2.20,1.32);
  \node[above] at (1.65,0.14) {$4$};
  \draw (1.10,0.87) ++(-138:0.26) -- ++(42:0.52);
  \draw (2.20,1.32) ++(-138:0.25) -- ++(42:0.50);

  % Lower half
  \draw (1.10,-0.87) -- (1.10,-0.12) -- (2.20,-0.12) -- (2.20,-1.17);
  \node[right] at (1.12,-0.52) {$r_2$};
  \node[right] at (2.22,-0.68) {$r_1$};
  \draw (1.10,-0.87) ++(138:0.25) -- ++(-42:0.50);
  \draw (2.20,-1.17) ++(138:0.26) -- ++(-42:0.52);
\end{tikzpicture}

\end{center}
\end{problem}
\clearpage


\begin{problem}{5.2.10}
用什么样的动系可把刚体做定点转动时用的 $\dot{\theta}, \dot{\varphi}, \dot{\psi}$ 三个角速度中的一个作为相对角速度，而另两个作为牵连角速度。
\end{problem}
\clearpage

